% !TEX root = ../main.tex
\section{Archived Extended Baselines}
\label{ch:appendix-archive}

Additional PageRank variants were scored on the same matched cohort and proxy families. They are not part of the primary EndorseRank--AWP research questions.

\dissertationSubheading[sec:seven-method-archive]{Seven-method alignment matrix}

Seven-method matrices remain in the evaluation archive (LP-PR, CW-AWP, LF-PR, and RiskProp-style variants). They are not drawn here. EndorseRank remains the allowance-family leader among social methods; AWP (and closely related transfer variants) lead on transfer.

The archive matrices show that:

\begin{itemize}
\item Methods built on transfer-like edges cluster near AWP on transfer and Sybil families.
\item Methods with very small edge counts can produce unstable or degenerate rankings; their runtimes may be low simply because $m$ is tiny.
\item RiskProp-style rows that mirror AWP numerically indicate shared transfer structure, not independent discovery.
\end{itemize}

\dissertationSubheading[sec:six-aave-archive]{Six-method Aave-oriented matrix}

An Aave-oriented preset remains in the evaluation archive. Several lending-event graphs are sparse on Arbitrum within the study window; delegation-based rows may be degenerate when \texttt{BorrowAllowanceDelegated} events were not extracted (Appendix~\ref{ch:appendix-eval}).

The Aave-oriented preset remains in the evaluation archive and is not drawn here.

These matrices motivate future work on credit-delegation edges and cross-protocol lending labels, but they do not alter Claims C1--C4.

\dissertationSubheading[sec:why-archived]{Why the comparison uses two methods}

The primary comparison is EndorseRank versus AWP for three reasons:

\begin{enumerate}
\item The novel construct is allowance-based endorsement. Comparing it to the established transfer baseline~\cend{3} isolates the research question.
\item Additional variants that share transfer edges with AWP add rows without adding independent constructs.
\item The outcome-native GF-PR reference used in the research proposal is built from the same GMX events as the trading-success proxies, so any alignment between them is construct overlap, not validation (Section~\ref{sub:gf-pr}); its three-method matrix is archived alongside the seven-method and six-Aave presets.
\end{enumerate}

\dissertationSubheading[sec:archive-repro]{Reproducing archive tables}

Archive tables are regenerated by the evaluation pipeline's table export with seven-method and six-Aave presets. CSV copies are recorded alongside the table fragments.

\dissertationSubheading[sec:archive-method-notes]{Notes on archived method identifiers}

The following identifiers appear in archive tables and CSV exports.

\begin{description}
\item[LP-PR] A transfer-graph variant with alternative loss-penalized or liquidity-oriented weighting. Empirically it tracks AWP closely on transfer families, which points to shared payment-flow structure.
\item[CW-AWP] A counterparty-weighted AWP variant that down-weights repeated transfers between the same pair. Slightly lowers transfer $\tau$ vs.\ AWP; same construct family.
\item[LF-PR] A low-frequency or lightly filtered graph; very small edge counts, near-zero runtimes. Rankings not comparable to full AWP: the sparse graph is a failed baseline.
\item[RiskProp] A risk-propagation style ranking on transfer-like edges, inspired by Lin et al.~\cend{12}. On this cohort it numerically mirrors AWP on several families, which implies the risk signal implemented did not diverge strongly from transfer centrality under the study's parameterization.
\item[LiqCall-PR / Borrow-PR / Repay-PR] Aave-oriented wallet-to-wallet graphs based on liquidation calls, borrows, and repays (where applicable). Edge counts on Arbitrum in the study window are small (tens to hundreds), so alignment is weak and runtimes are negligible.
\item[Delegation-PR] Intended to leverage credit-delegation allowances; with zero delegation edges extracted, scores are degenerate and $\tau$ is not interpretable.
\end{description}

Archive rows are not peer-reviewed baselines.

\dissertationSubheading[sec:archive-decision-log]{Criteria for the two-method comparison}

Chapters~\ref{ch:results}--\ref{ch:discussion} focus on EndorseRank and AWP. The other variants do not meet the following criteria:

\begin{enumerate}
\item Construct independence: a baseline should implement a different edge semantics, rather than simply being a reweighting of transfers.
\item Data adequacy: the edge counts should be large enough to support stable PageRank on the matched cohort.
\item Scope: the contribution is allowance-based endorsement, not a survey of all possible PageRank variants.
\end{enumerate}

Based on these criteria, EndorseRank (allowance) and AWP (transfer) form a minimal sufficient set. The seven-method and six-Aave matrices remain in the archive; they are not included in Claims~C1--C4.

